% GENERATED FILE. Edit canonical lessons, then run npm run generate:books.
% canonical-source-hash: fnv1a64:a4e06422ad3e31fd
% canonical-lessons: JA-C139-desuka, JA-C139-chigaimasu, JA-C139-janai, JA-C139-nanimo, JA-C139-daremo, JA-C139-dokonimo

\chapter{Is It? No: Nothing, Nobody, Nowhere}
\label{ch:ja-is-it-no-nothing-nobody-nowhere}
\hlchaptermodality{139}

\begin{chapteropening}
\textbf{By the end of this chapter:} \emph{I can ask a yes-or-no question with ka, answer that something is not so, and say nothing, nobody and nowhere with a negative verb.}

Ask a yes-or-no question with desu ka, answer no with chigaimasu, say what is not so with ja nai, and say nothing, nobody and nowhere, each with a verb in -masen.
\end{chapteropening}

\section[\texorpdfstring{\ja{ですか} (desu ka)}{ですか (desu ka)}]{\ja{ですか} (desu ka) — is it ...? (the polite yes-or-no question)}

\label{lesson:JA-C139-desuka}



\begin{quote}
Four recalls before the new word.

\begin{itemize}
\raggedright
  \item \emph{Recall:} say \emph{I go to bed}, politely — \textbf{R1}, one lesson back
  \item \emph{Recall:} say \emph{I get up}, politely — \textbf{R2}, five lessons back
  \item \emph{Recall:} say \emph{a wind}, or \emph{a cold} — \textbf{R3}, twenty lessons back
  \item \emph{Recall:} say \emph{tea} — \textbf{R4}, eighty lessons back
\end{itemize}
\end{quote}

\subsection*{You'll want to know: \ja{ですか}}
\textbf{\ja{ですか}} — \emph{desu ka} — \textquotedblleft{}is it ...?\textquotedblright{}. \textbf{\ja{です}}, \emph{desu}, says politely that something is so. Put \textbf{\ja{か}}, \emph{ka}, at the end, and the sentence becomes a question that can be answered yes or no. The order of the words does not change, and the voice rises a little on \emph{ka}.

\begin{quote}
\textbf{\ja{これは} \ja{みずですか。}} — \emph{kore wa mizu desu ka} — \textquotedblleft{}Is this water?\textquotedblright{} \textbf{\ja{はい、みずです。}} — \emph{hai, mizu desu} — \textquotedblleft{}Yes, it is water.\textquotedblright{}
\end{quote}

A verb takes \textbf{\ja{か}} the same way: \textbf{\ja{ます}} becomes \textbf{\ja{ますか}}.

\begin{quote}
\textbf{\ja{あした} \ja{はたらきますか。}} — \emph{ashita hatarakimasu ka} — \textquotedblleft{}Are you working tomorrow?\textquotedblright{}
\end{quote}

\subsection*{Your turn}
\begin{itemize}
\raggedright
  \item \emph{Say it:} \emph{desu ka}
  \item \emph{Say it:} \emph{kore wa mizu desu ka} — Is this water?
  \item \emph{Recall:} turn \emph{I go home} into a question: \emph{kaerimasu ka}
\end{itemize}

\begin{tcolorbox}[breakable,colback=teal!4,colframe=teal!35!black,title={Before you move on}]
How do you ask \textquotedblleft{}is it ...?\textquotedblright{}, politely? (\textbf{\ja{ですか}}, \emph{desu ka}.) Say it once more.
\end{tcolorbox}
\section[\texorpdfstring{\ja{ちがいます} (chigaimasu)}{ちがいます (chigaimasu)}]{\ja{ちがいます} (chigaimasu) — no, that is not right; it is different (polite)}

\label{lesson:JA-C139-chigaimasu}



\begin{quote}
Four recalls before the new word.

\begin{itemize}
\raggedright
  \item \emph{Recall:} ask \emph{is it ...?}, politely — \textbf{R1}, one lesson back
  \item \emph{Recall:} say \emph{I work}, politely — \textbf{R2}, five lessons back
  \item \emph{Write it:} \textbf{\ja{ぜ}} from memory — \textbf{R3}, twenty lessons back
  \item \emph{Write it:} small \textbf{\ja{ゃ}} from memory — \textbf{R4}, eighty lessons back
\end{itemize}
\end{quote}

\subsection*{You'll want to know: \ja{ちがいます}}
\textbf{\ja{ちがいます}} — \emph{chigaimasu} — \textquotedblleft{}it is different\textquotedblright{}, and so, as an answer, \textquotedblleft{}no, that is not right\textquotedblright{}. It is the polite form of the verb \textbf{\ja{ちがう}}, \emph{chigau}, \textquotedblleft{}to differ\textquotedblright{}. When someone asks whether a thing is so and it is not, this is the everyday answer.

\begin{quote}
\textbf{\ja{これは} \ja{おちゃですか。}} — \emph{kore wa ocha desu ka} — \textquotedblleft{}Is this tea?\textquotedblright{} \textbf{\ja{いいえ、ちがいます。みずです。}} — \emph{iie, chigaimasu. mizu desu} — \textquotedblleft{}No, it isn't. It is water.\textquotedblright{}
\end{quote}

\subsection*{Your turn}
\begin{itemize}
\raggedright
  \item \emph{Say it:} \emph{chigaimasu}
  \item \emph{Say it:} \emph{iie, chigaimasu} — No, that is not right
  \item \emph{Recall:} ask \emph{Is this tea?}, then answer \emph{No, it isn't}
\end{itemize}

\begin{tcolorbox}[breakable,colback=teal!4,colframe=teal!35!black,title={Before you move on}]
How do you answer \textquotedblleft{}no, that is not right\textquotedblright{}? (\textbf{\ja{ちがいます}}, \emph{chigaimasu}.) Say it once more.
\end{tcolorbox}
\section[\texorpdfstring{\ja{じゃない} (ja nai)}{じゃない (ja nai)}]{\ja{じゃない} (ja nai) — is not, is not so (plain)}

\label{lesson:JA-C139-janai}



\begin{quote}
Three recalls before the new word.

\begin{itemize}
\raggedright
  \item \emph{Recall:} answer \emph{no, that is not right} — \textbf{R1}, one lesson back
  \item \emph{Recall:} say \emph{I walk}, politely — \textbf{R2}, five lessons back
  \item \emph{Recall:} say \emph{a doctor} — \textbf{R4}, eighty lessons back
\end{itemize}
\end{quote}

\subsection*{You'll want to know: \ja{じゃない}}
\textbf{\ja{じゃない}} — \emph{ja nai} — \textquotedblleft{}is not\textquotedblright{}. It is the plain way, among friends, to say that something is not so. The polite way is \textbf{\ja{じゃありません}}, \emph{ja arimasen}.

\begin{quote}
\textbf{\ja{みずじゃない。}} — \emph{mizu ja nai} — \textquotedblleft{}It isn't water.\textquotedblright{}
\end{quote}

Verbs say \emph{not} with their ending. In the polite form, \textbf{\ja{ます}} becomes \textbf{\ja{ません}}, \emph{masen}:

\begin{quote}
\textbf{\ja{はたらきます}} $\to$ \textbf{\ja{はたらきません}} — \emph{hatarakimasen} — \textquotedblleft{}I don't work\textquotedblright{}
\end{quote}

\begin{quote}
\textbf{\ja{ねます}} $\to$ \textbf{\ja{ねません}} — \emph{nemasen} — \textquotedblleft{}I don't go to bed\textquotedblright{}
\end{quote}

\subsection*{Your turn}
\begin{itemize}
\raggedright
  \item \emph{Say it:} \emph{ja nai}
  \item \emph{Say it:} \emph{mizu ja nai} — It isn't water
  \item \emph{Recall:} say \emph{I get up}, then \emph{I don't get up}: \emph{okimasu}, \emph{okimasen}
\end{itemize}

\begin{tcolorbox}[breakable,colback=teal!4,colframe=teal!35!black,title={Before you move on}]
How do you say \textquotedblleft{}is not\textquotedblright{}, plainly? (\textbf{\ja{じゃない}}, \emph{ja nai}.) Say it once more.
\end{tcolorbox}
\section[\texorpdfstring{\ja{なにも} (nani mo)}{なにも (nani mo)}]{\ja{なにも} (nani mo) — nothing, not anything (with a negative verb)}

\label{lesson:JA-C139-nanimo}



\begin{quote}
Four recalls before the new word.

\begin{itemize}
\raggedright
  \item \emph{Recall:} say \emph{is not}, plainly — \textbf{R1}, one lesson back
  \item \emph{Recall:} say \emph{I go home}, politely — \textbf{R2}, five lessons back
  \item \emph{Recall:} say \emph{full; a lot} — \textbf{R3}, twenty lessons back
  \item \emph{Recall:} say \emph{a little; a moment} — \textbf{R4}, eighty lessons back
\end{itemize}
\end{quote}

\subsection*{You'll want to know: \ja{なにも}}
\textbf{\ja{なにも}} — \emph{nani mo} — \textquotedblleft{}nothing\textquotedblright{}, \textquotedblleft{}not anything\textquotedblright{}. It pairs with a negative verb: Japanese says \textquotedblleft{}anything ... not\textquotedblright{}, where English says \textquotedblleft{}nothing\textquotedblright{}.

\begin{quote}
\textbf{\ja{なにも} \ja{のみません。}} — \emph{nani mo nomimasen} — \textquotedblleft{}I don't drink anything.\textquotedblright{}
\end{quote}

\textbf{\ja{のむ}}, \emph{nomu}, \textquotedblleft{}to drink\textquotedblright{}, gives \textbf{\ja{のみます}} \emph{nomimasu} and then \textbf{\ja{のみません}} \emph{nomimasen}. With \textbf{\ja{ます}} instead of \textbf{\ja{ません}} the sentence does not work: \textbf{\ja{なにも}} needs the \emph{not}.

\subsection*{Your turn}
\begin{itemize}
\raggedright
  \item \emph{Say it:} \emph{nani mo}
  \item \emph{Say it:} \emph{nani mo nomimasen} — I don't drink anything
  \item \emph{Recall:} say \emph{I don't eat anything}: \emph{nani mo tabemasen}
\end{itemize}

\begin{tcolorbox}[breakable,colback=teal!4,colframe=teal!35!black,title={Before you move on}]
How do you say \textquotedblleft{}nothing\textquotedblright{}, with a negative verb? (\textbf{\ja{なにも}}, \emph{nani mo}.) Say it once more.
\end{tcolorbox}
\section[\texorpdfstring{\ja{だれも} (dare mo)}{だれも (dare mo)}]{\ja{だれも} (dare mo) — nobody, not anyone (with a negative verb)}

\label{lesson:JA-C139-daremo}



\begin{quote}
Four recalls before the new word.

\begin{itemize}
\raggedright
  \item \emph{Recall:} say \emph{nothing}, with a negative verb — \textbf{R1}, one lesson back
  \item \emph{Recall:} say \emph{I go to bed}, politely — \textbf{R2}, five lessons back
  \item \emph{Write it:} \textbf{\ja{ぱ}} from memory — \textbf{R3}, twenty lessons back
  \item \emph{Write it:} small \textbf{\ja{ょ}} from memory — \textbf{R4}, eighty lessons back
\end{itemize}
\end{quote}

\subsection*{You'll want to know: \ja{だれも}}
\textbf{\ja{だれも}} — \emph{dare mo} — \textquotedblleft{}nobody\textquotedblright{}, \textquotedblleft{}not anyone\textquotedblright{}. Like \textbf{\ja{なにも}}, it needs a negative verb.

\begin{quote}
\textbf{\ja{ここには} \ja{だれも} \ja{いません。}} — \emph{koko ni wa dare mo imasen} — \textquotedblleft{}There is nobody here.\textquotedblright{}
\end{quote}

\textbf{\ja{いません}}, \emph{imasen}, is the negative of \textbf{\ja{います}}, \emph{imasu}, \textquotedblleft{}someone is there\textquotedblright{}, the verb Japanese uses for people and animals.

\subsection*{Your turn}
\begin{itemize}
\raggedright
  \item \emph{Say it:} \emph{dare mo}
  \item \emph{Say it:} \emph{koko ni wa dare mo imasen} — There is nobody here
  \item \emph{Recall:} say \emph{nothing}, then \emph{nobody}, each with a negative verb
\end{itemize}

\begin{tcolorbox}[breakable,colback=teal!4,colframe=teal!35!black,title={Before you move on}]
How do you say \textquotedblleft{}nobody\textquotedblright{}, with a negative verb? (\textbf{\ja{だれも}}, \emph{dare mo}.) Say it once more.
\end{tcolorbox}
\section[\texorpdfstring{\ja{どこにも} (doko ni mo)}{どこにも (doko ni mo)}]{\ja{どこにも} (doko ni mo) — nowhere, not anywhere (with a negative verb)}

\label{lesson:JA-C139-dokonimo}



\begin{quote}
Four recalls before the new word.

\begin{itemize}
\raggedright
  \item \emph{Recall:} say \emph{nobody}, with a negative verb — \textbf{R1}, one lesson back
  \item \emph{Recall:} ask \emph{is it ...?}, politely — \textbf{R2}, five lessons back
  \item \emph{Recall:} say \emph{worry} — \textbf{R3}, twenty lessons back
  \item \emph{Recall:} say \emph{a library} — \textbf{R4}, eighty lessons back
\end{itemize}
\end{quote}

\subsection*{You'll want to know: \ja{どこにも}}
\textbf{\ja{どこにも}} — \emph{doko ni mo} — \textquotedblleft{}nowhere\textquotedblright{}, \textquotedblleft{}not anywhere\textquotedblright{}. It is \textbf{\ja{どこ}}, \textquotedblleft{}where\textquotedblright{}, with \textbf{\ja{に}} for the place you go to and \textbf{\ja{も}}, and like the other two it needs a negative verb.

\begin{quote}
\textbf{\ja{あしたは} \ja{どこにも} \ja{いきません。}} — \emph{ashita wa doko ni mo ikimasen} — \textquotedblleft{}I'm not going anywhere tomorrow.\textquotedblright{}
\end{quote}

Now you can ask and answer:

\begin{quote}
\textbf{\ja{あした} \ja{はたらきますか。}} — \textquotedblleft{}Are you working tomorrow?\textquotedblright{} \textbf{\ja{いいえ、はたらきません。どこにも} \ja{いきません。}} — \textquotedblleft{}No, I'm not. I'm not going anywhere.\textquotedblright{}
\end{quote}

\subsection*{Your turn}
\begin{itemize}
\raggedright
  \item \emph{Say it:} \emph{doko ni mo}
  \item \emph{Say it:} \emph{ashita wa doko ni mo ikimasen} — I'm not going anywhere tomorrow
  \item \emph{Recall:} ask \emph{Is this water?} (\emph{desu ka}), answer \emph{No, that is not right} (\emph{chigaimasu}), say \emph{It isn't water} plainly (\emph{ja nai}), then \emph{nothing}, \emph{nobody} and \emph{nowhere}, each with a negative verb
\end{itemize}

\begin{tcolorbox}[breakable,colback=teal!4,colframe=teal!35!black,title={Before you move on}]
Is it ...? (\textbf{\ja{ですか}}, \emph{desu ka}.) No, that is not right? (\textbf{\ja{ちがいます}}, \emph{chigaimasu}.) Is not, plainly? (\textbf{\ja{じゃない}}, \emph{ja nai}.) Nothing? (\textbf{\ja{なにも}}, \emph{nani mo}.) Nobody? (\textbf{\ja{だれも}}, \emph{dare mo}.) Nowhere? (\textbf{\ja{どこにも}}, \emph{doko ni mo}.)
\end{tcolorbox}
